\chapter{Descripción del Problema}

Las Amenazas Persistentes Avanzadas (APT) constituyen una de las formas de ataque más difíciles de enfrentar en la ciberseguridad actual. A diferencia de un ataque puntual, una APT corresponde a una campaña dirigida y sigilosa que puede extenderse durante meses y que se desarrolla en múltiples etapas sucesivas, desde el reconocimiento inicial del objetivo hasta la exfiltración de información \parencite{hutchins2011killchain}. Esta naturaleza prolongada y por etapas es precisamente lo que les permite operar sin ser detectadas: cada acción individual del atacante, observada de forma aislada, resulta indistinguible de la actividad legítima de la red.

El problema radica en que los mecanismos de detección tradicionales no están diseñados para capturar este comportamiento. Las herramientas basadas en firmas estáticas, así como buena parte de los sistemas de detección de intrusiones convencionales, evalúan eventos de manera aislada y no la secuencia temporal completa que caracteriza a una campaña de este tipo, por lo que las APT logran evadirlos con frecuencia \parencite{alshamrani2019survey}. A esta limitación se suma una segunda dificultad, esta vez en el plano de los datos disponibles para investigación: los conjuntos de datos de intrusiones genéricos presentan una diferencia demasiado evidente entre el tráfico normal y el anómalo, lo que no representa la sofisticación real de un atacante APT, y su distribución de clases no refleja las condiciones reales de operación \parencite{myneni2020dapt2020}. En escenarios reales, la actividad maliciosa constituye una fracción mínima del tráfico total, un desbalance severo que dificulta tanto el entrenamiento como la evaluación de modelos de detección.

Es necesario distinguir entre el problema general y la brecha específica que aborda esta investigación. La detección de APT ha sido ampliamente estudiada, y existen enfoques que incorporan la dimensión temporal del ataque mediante modelos de progresión por etapas. Por tanto, la brecha está en las condiciones bajo las cuales estas propuestas han sido validadas. Los trabajos existentes tienden a evaluarse sobre conjuntos de datos con distribuciones de clase artificialmente equilibradas o con una separación demasiado evidente entre tráfico legítimo y malicioso, escenarios que no reproducen las condiciones reales de operación. En consecuencia, permanece abierta la pregunta de si los enfoques probabilísticos secuenciales mantienen su desempeño cuando la actividad maliciosa constituye una fracción mínima del tráfico total, condición que caracteriza a los entornos reales y que este trabajo aborda de manera explícita.

Las consecuencias de esta brecha son concretas. Una campaña que permanece activa durante meses sin ser detectada implica pérdida sostenida de información confidencial, un tiempo de respuesta que se mide en semanas o meses en lugar de horas, y la imposibilidad de intervenir en las etapas tempranas del ataque, cuando el daño aún es limitado. Existe, por tanto, la necesidad de enfoques de detección capaces de razonar sobre la progresión temporal de un ataque y de hacerlo bajo las condiciones de desbalance de datos que caracterizan a los escenarios reales. Esta situación se evidencia en trabajos recientes: al evaluar modelos de detección de anomalías ampliamente utilizados sobre un escenario que emula seis semanas de actividad organizacional real, la actividad correspondiente a la campaña APT resulta prácticamente indetectable \parencite{myneni2023unraveled}.